\section{Giao thức Mạng và Mô hình I/O}

\subsection[Mô hình Client-Server và I/O Multiplexing]{Mô hình Client-Server và I/O Multiplexing (\texttt{epoll})}

\begin{enumerate}
    \item Cơ sở lý thuyết \\
% REMOVED BLANK LINE
    Mô hình Thread-per-connection phát sinh chi phí chuyển ngữ cảnh (Context Switching) lớn đối với bài toán C10K. Do đó, kỹ thuật I/O Multiplexing với API \texttt{epoll} (Linux) được áp dụng để một luồng có thể giám sát nhiều mô tả tệp (File Descriptor - FD). Thông qua cấu trúc cây đỏ đen (Red-Black Tree) và danh sách liên kết (Linked List), \texttt{epoll} đạt độ phức tạp $\mathcal{O}(1)$ trong quá trình xử lý sự kiện.

    \item Triển khai hệ thống \\
% REMOVED BLANK LINE
    Thành phần \texttt{TcpServer} sử dụng I/O phi đồng bộ (Non-blocking I/O) và \texttt{epoll} để điều phối kết nối TCP theo quy trình sau:
    \begin{itemize}[label={-}]
        \item Khởi tạo điểm cuối: Chuyển trạng thái Socket lắng nghe sang phi đồng bộ thông qua \texttt{fcntl}.
        \item Khởi tạo I/O Multiplexer: Tạo đối tượng \texttt{epoll} (\texttt{epoll\_create1}) và đăng ký Socket lắng nghe với cờ \texttt{EPOLLIN} (\texttt{epoll\_ctl}).
        \item Vòng lặp sự kiện: Thực thi \texttt{epoll\_wait} để thu thập sự kiện I/O.
        \begin{itemize}[label={$\circ$}]
            \item Sự kiện kết nối: Gọi \texttt{accept}, chuyển \texttt{clientFd} sang trạng thái phi đồng bộ, đăng ký vào \texttt{epoll} (cờ \texttt{EPOLLIN | EPOLLRDHUP}) và khởi tạo \texttt{ClientSession}.
            \item Sự kiện dữ liệu: \texttt{TcpServer} tiếp nhận byte thô (\texttt{recv}) và giao cho \texttt{ClientSession} xử lý.
        \end{itemize}
    \end{itemize}
\end{enumerate}

\subsection[Thiết kế Định dạng Gói tin]{Thiết kế Định dạng Gói tin (Packet Format)}

\begin{enumerate}
    \item Cơ sở lý thuyết \\
% REMOVED BLANK LINE
    TCP là giao thức hướng luồng (Stream-oriented), gây ra hiện tượng dính gói (Sticky Packet) và phân mảnh (Half Packet). Tầng ứng dụng sử dụng giao thức nhị phân (Binary Protocol) để xác định ranh giới gói tin, tối ưu băng thông và giảm chi phí phân tích cú pháp (Parsing Overhead).

    \item Cấu trúc gói tin \\
% REMOVED BLANK LINE
    Hệ thống thiết kế giao thức nhị phân gồm Header (17 Bytes) và Payload, được đặc tả chi tiết tại Bảng \ref{tab:packet_format}.
\end{enumerate}


\renewcommand{\arraystretch}{1.4}
    \begin{longtable}{
        >{\raggedright\arraybackslash}p{3cm}
        >{\raggedright\arraybackslash}p{4cm}
        >{\raggedright\arraybackslash}p{2.5cm}
        >{\raggedright\arraybackslash}p{5.5cm}
    }
        \caption{Cấu trúc Gói tin Nhị phân} \label{tab:packet_format} \\
        \toprule
        \textbf{Phân đoạn} & \textbf{Trường dữ liệu} & \textbf{Kích thước} & \textbf{Mô tả chi tiết} \\
        \midrule
        \endfirsthead
        
        \caption[]{Cấu trúc Gói tin Nhị phân (tiếp theo)} \\
        \toprule
        \textbf{Phân đoạn} & \textbf{Trường dữ liệu} & \textbf{Kích thước} & \textbf{Mô tả chi tiết} \\
        \midrule
        \endhead
        
        \bottomrule
        \endfoot
        
        \bottomrule
        \endlastfoot

        Header & \texttt{magic} (Nhận diện) & 2 Bytes & Chuỗi byte nhận diện tính hợp lệ. \\
        & \texttt{version} (Phiên bản) & 1 Byte & Định danh phiên bản giao thức. \\
        & \texttt{msg\_type} (Loại tin) & 1 Byte & Định nghĩa nghiệp vụ (đăng nhập, trò chuyện...). \\
        & \texttt{flags} (Cờ điều khiển) & 1 Byte & Trạng thái gói tin (mã hóa, nén). \\
        & \texttt{sequence\_num} (Thứ tự) & 4 Bytes & Hỗ trợ ráp mảnh. \\
        & \texttt{payload\_length} (Kích thước) & 4 Bytes & Giải quyết hiện tượng Sticky Packet. \\
        & \texttt{checksum} (Mã kiểm tra) & 4 Bytes & Xác thực toàn vẹn dữ liệu (Integrity). \\
        \midrule
        
        Payload & Dữ liệu JSON & Tùy biến & Đóng gói tham số nghiệp vụ (Mã hóa AES-256). \\
    \end{longtable}
    

Chu trình phân giải gói tin (Packet Parsing) gồm các bước:
\begin{itemize}[label={-}]
    \item Đóng gói: Tuần tự hóa JSON, mã hóa và gắn Header (17 Bytes) trước Payload.
    \item Phân tích Header: \texttt{ClientSession} đọc 17 bytes đầu tiên để trích xuất các tham số điều khiển (\texttt{msg\_type}, \texttt{payload\_length}, \texttt{magic}, \texttt{checksum}).
    \item Đọc Payload: Cấp phát bộ nhớ dựa trên \texttt{payload\_length} để nhận luồng Payload và xử lý hiện tượng phân mảnh.
    \item Giải mã và định tuyến: \texttt{CryptoEngine} giải mã khối dữ liệu, sau đó hệ thống định tuyến theo \texttt{msg\_type} đến bộ xử lý tương ứng.
\end{itemize}